\section{Background and Related Work}
\label{sec:related}

\subsection{Soft Bindings, Provenance Standards and Regulation}
The C2PA Technical Specification 2.4 \cite{c2pa2026spec} binds a manifest to its asset with a cryptographic \emph{hard binding} and optionally with one or more \emph{soft bindings}: content identifiers that are either not statistically unique, such as a fingerprint, or embedded as an invisible watermark (Sec.~2.3.13). A fingerprint is defined as ``a set of inherent properties computable from digital content that identifies the content or near duplicates of it'' (Sec.~2.4.2). A soft binding is recorded in a \texttt{c2pa.soft-binding} assertion carrying an algorithm identifier, the value and optional scope (Sec.~18.10), and it lets a validator that holds content without a manifest query a repository through the Soft Binding Resolution API \cite{c2pa2026softbindingapi}. Version 2.4 states that the algorithm \emph{should} be one named in the C2PA soft-binding algorithm list (Sec.~9.3.2) and that a soft binding shall not be used as a hard binding (Sec.~9.3.1). At the snapshot we used (commit \texttt{5c864efb}, 16 September 2026) the list held \val{fp:c2pa/list_total} entries, \val{fp:c2pa/list_watermark} watermarks and \val{fp:c2pa/list_fingerprint} fingerprints \cite{c2pasoftbindinglist}. Of the fingerprints, only ISCC (\texttt{io.iscc.v0}) is an open, multi-modal standard \cite{iso24138}; the others are proprietary services or, in one case, a cryptographic digest registered as a soft binding. Collomosse and Parsons describe the combination of manifest, watermark and fingerprint as the three pillars of durable provenance \cite{collomosse2024durable}, and Adobe's deep image comparator \cite{black2021dic} and SImProv \cite{black2022simprov} combine retrieval with localization of the edited region; neither is released as code.

Article~50(2) of the EU Artificial Intelligence Act requires providers of systems that generate synthetic audio, image, video or text to ensure that outputs are marked in a machine-readable format and detectable as artificially generated or manipulated. The technical solutions must be effective, interoperable, robust and reliable as far as technically feasible \cite{eu2024aiact}. The Digital Omnibus on AI, in force since 27 July 2026, postponed this marking obligation to 2 December 2026 for systems placed on the market before 2 August 2026 \cite{eu2026omnibus}. The Act does not name a technique or set thresholds. Imperceptibility, survival under processing, false-positive rate and cost are therefore left to evidence of the kind a comparative benchmark provides.

\subsection{Invisible Watermarking}
\paragraph{Images.} Cox et al.\ introduced spread-spectrum watermarking for multimedia, spreading a pseudo-random sequence over perceptually significant transform coefficients so that the mark cannot be removed without degrading the host \cite{cox1997secure}. Later transform-domain schemes embedded the mark in the DCT domain with a detector that does not need the original image \cite{barni1998dct}, or combined the discrete wavelet transform with singular value decomposition \cite{ganic2004dwtsvd}. The DWT-DCT and DWT-DCT-SVD encoders of the invisible-watermark library \cite{invisiblewatermark} are open implementations in this family. Their documentation warns that decoding is not guaranteed even on unmodified images.

Deep watermarking replaced hand-designed transforms with learned encoders and decoders. HiDDeN trained both networks jointly and showed that the encodings can be made robust to noise applied between them \cite{zhu2018hidden}. StegaStamp extended the idea to printed and re-photographed images \cite{tancik2020stegastamp}, and MBRS handled non-differentiable JPEG compression by mixing real and simulated compression within each mini-batch \cite{jia2021mbrs}. RivaGAN added an attention-based embedder and two adversarial critics, one for visual quality and one for robustness, and was designed for video \cite{zhang2019rivagan}. RoSteALS moved embedding into the latent space of a frozen pretrained autoencoder \cite{bui2023rosteals}. TrustMark, a GAN-based method with spatio-spectral losses, targets images of arbitrary resolution and adds a network that removes its own watermark for re-watermarking \cite{bui2025trustmark}. InvisMark was designed for high-resolution generated images and reports a PSNR near 51\,dB, an SSIM near 0.998 and 256-bit payloads \cite{xu2025invismark}. The Watermark Anything Model (WAM) segments an image into watermarked and unwatermarked regions and recovers several 32-bit messages from small areas \cite{sander2025wam}. Pixel Seal replaces pixel-wise perceptual losses with adversarial-only training, stabilizes it with a three-stage schedule, and applies just-noticeable-difference attenuation at high resolution \cite{soucek2025pixelseal}. Petrov et al.\ derived upper bounds on watermark capacity under PSNR and linear robustness constraints and trained ChunkySeal, a scaled version of Video Seal that carries 1024 bits \cite{petrov2025chunkyseal}.

A separate line embeds the mark during generation. Stable Signature, for example, fine-tunes the latent decoder of a diffusion model, conditioned on a binary signature, so that every generated image carries it \cite{fernandez2023stable}. Such methods cannot mark existing content, so our study covers post-hoc methods only. Each of the papers above evaluates its method against baselines of its own choosing, on its own images, at its own embedding strength and with its own selection of metrics, mostly PSNR and SSIM. Their reported numbers are therefore not directly comparable.

\paragraph{Video.} Hartung and Girod applied spread-spectrum embedding to both uncompressed and compressed video \cite{hartung1998video}. Among learned methods, DVMark uses a multiscale network for video watermarking \cite{luo2025dvmark}, and RivaGAN operates on frame sequences \cite{zhang2019rivagan}. Video Seal trains an embedder and an extractor with video codecs among the training transformations. It also propagates the watermark temporally, so that an image model can be applied without watermarking every high-resolution frame \cite{fernandez2024videoseal}. Pixel Seal adapts to video by pooling the watermark over time \cite{soucek2025pixelseal}. In deployment, video is usually re-encoded with H.264/AVC \cite{wiegand2003h264} or H.265/HEVC \cite{sullivan2012hevc}, and inter-frame prediction together with coarse quantization can attenuate low-energy residuals. Few video watermarking models have openly licensed weights, and independent evidence is thin: VideoMarkBench, the dedicated video benchmark discussed below, uses clips produced by video generators rather than camera footage \cite{jiang2025videomarkbench}.

\paragraph{Audio.} Bender et al.\ surveyed early data-hiding techniques for images and audio \cite{bender1996techniques}, and Kirovski and Malvar adapted spread-spectrum watermarking to audio signals \cite{kirovski2003spread}. WavMark encodes up to 32 bits in each second of audio and reports a mean bit error rate of 0.48\% over ten attacks on 10-20\,s clips \cite{chen2023wavmark}. AudioSeal trains a generator and a detector with a localization loss for sample-level detection and uses a perceptual loss inspired by auditory masking; its single-pass detector is intended for large-scale screening \cite{sanroman2024audioseal}. SilentCipher adds psychoacoustic thresholding and pseudo-differentiable compression layers and operates on 44.1\,kHz audio \cite{singh2024silentcipher}. Perth is an open library whose main watermarker is a neural implicit model. Its documentation claims survival of compression and resampling, but no peer-reviewed evaluation is available \cite{perth2025}. WavMark and AudioSeal were motivated by voice cloning and evaluated mainly on speech, and the four methods use different native sample rates, so a shared protocol has to fix resampling and codec settings before their numbers can be compared.

\subsection{Content Fingerprinting}
\paragraph{Perceptual hashes.} Classical perceptual hashes reduce an image to a short binary code that changes little under common processing. The average and difference hashes compare block means \cite{krawetz2011looks, krawetz2013kind}; pHash thresholds low-frequency DCT coefficients at their median \cite{zauner2010phash}; the Marr--Hildreth and block-mean hashes use edge and block statistics \cite{zauner2010phash, yang2006blockhash}, and implementations are collected in imagehash \cite{buchner_imagehash}, Blockhash \cite{commonsmachinery_blockhash} and OpenCV's \texttt{img\_hash} module \cite{opencv_imghash}. Meta's PDQ applies a DCT to a filtered 64$\times$64 luminance image and keeps 256 bits together with a quality score \cite{meta2019pdq}; it is deployed in industry hash-sharing programmes. The ISCC Image-Code of ISO 24138 is a DCT hash of a 32$\times$32 greyscale image after trimming uniform borders \cite{iso24138}. The ISCC benchmark TwinSpect \cite{iscc_twinspect} evaluates such compact codes, but to our knowledge no study has compared them with learned descriptors under a common false-positive calibration.

\paragraph{Learned copy-detection descriptors.} The Image Similarity Challenge 2021 and its DISC21 dataset established copy detection as a retrieval problem with a micro-average precision metric over one million references \cite{douze2021isc}. SSCD trains a descriptor with self-supervised contrastive learning, entropy regularisation and heavy augmentation, and outperformed classification-oriented self-supervised features by a wide margin \cite{pizzi2022sscd}; the winning descriptor-track entry used a large memory bank and negative-embedding subtraction \cite{yokoo2021isc}. General-purpose embeddings such as DINOv2 \cite{oquab2024dinov2} and CLIP \cite{radford2021clip, cherti2023openclip} capture semantics rather than copy identity, which helps under heavy geometric change but risks matching different photographs of the same object. ISCC's experimental semantic image code binarises a descriptor derived from the ISC21 winner \cite{iscc_sci}, and DINOHash adversarially fine-tunes DINOv2 into a 96-bit hash \cite{singhi2025dinohash}. Local features with geometric verification, from SIFT \cite{lowe2004sift} and ORB \cite{rublee2011orb} with RANSAC \cite{fischler1981ransac} to DISK \cite{tyszkiewicz2020disk} with LightGlue \cite{lindenberger2023lightglue}, remain the standard way to confirm a partial copy and to register a query onto its source.

\paragraph{Video and audio fingerprints.} For video, the MPEG-7 video signature \cite{paschalakis2012mpeg7} underlies the ISCC Video-Code, Meta's TMK+PDQF pools per-frame PDQ features with a temporal match kernel \cite{poullot2015tmk, meta2019pdq}, and vPDQ matches clips by the share of frames with a close PDQ hash \cite{meta_vpdq}. The 2023 Video Similarity Challenge showed that frame-level learned descriptors with temporal alignment dominate for copied segments \cite{pizzi2024vsc}. Audio identification is mature: landmark hashing \cite{wang2003shazam} and sub-band energy fingerprints \cite{haitsma2002robust} underlie tools such as audfprint \cite{ellis_audfprint}, Olaf \cite{six2023olaf}, Panako \cite{six2022panako} and Chromaprint \cite{lalinsky_chromaprint}, on which the ISCC Audio-Code is built. Neural audio fingerprints learn segment embeddings with contrastive objectives \cite{chang2021nafp, araz2025nmfp}.

\subsection{Benchmarks, Attacks and Security}
Petitcolas et al.\ catalogued attacks on early copyright marking systems and argued that such systems need common evaluation \cite{petitcolas1998attacks}. Recent benchmarks were motivated by generative models. WAVES evaluates image watermarks on detection and identification under conventional distortions, regeneration by diffusion models and adversarial attacks \cite{an2024waves}. W-Bench tests eleven image watermarking methods against regeneration, global and local editing and image-to-video generation, and proposes VINE, a watermark built on a pretrained diffusion model \cite{lu2025vine}. MarkDiffusion is a toolkit for in-generation watermarking of latent diffusion models, with shared evaluation pipelines for detectability, robustness and output quality \cite{pan2025markdiffusion}. AudioMarkBench evaluates three audio watermarking methods against 15 perturbations on a multilingual speech dataset under no-box, black-box and white-box threat models \cite{liu2024audiomarkbench}. VideoMarkBench evaluates four video watermarking methods against 12 perturbations on clips produced by three video generators \cite{jiang2025videomarkbench}.

Attack papers bound what any benchmark can claim. Zhao et al.\ proved that regeneration attacks, which add noise and then reconstruct the image, remove pixel-level invisible watermarks \cite{zhao2024removable}. Jiang et al.\ showed that small adversarial perturbations evade watermark detection with less visible change than JPEG compression or blurring \cite{jiang2023evading}. Our watermarking track measures robustness to routine processing (codecs, resizing, cropping, filtering and resampling), that is, survival of ordinary distribution rather than security against a motivated adversary; the fingerprinting track adds white-box, transfer, forgery and inversion attacks (Section~\ref{sec:fp-res-security}).

Perceptual hashes were designed for benign processing, not for an adversary. Gradient attacks against Apple's NeuralHash produced collisions and evasions with imperceptible perturbations \cite{struppek2022neuralhash}; black-box attacks evaded PDQ and PhotoDNA and forged targeted collisions \cite{jain2022adversarial, prokos2023squint, dolhansky2020collision}; and inversion attacks reconstruct recognisable images from PDQ, PhotoDNA and NeuralHash codes \cite{hawkes2025inversion, madden2024robustness}. The last result matters for provenance: a soft-binding value is written into the manifest, so a descriptor that can be inverted discloses content to anyone who reads the manifest.

Following standard practice for binomial rates and resampled statistics, we report exact Clopper-Pearson intervals \cite{clopper1934} and bootstrap intervals \cite{efron1979bootstrap}.

\subsection{Watermarks and Fingerprints Together}
Invisible watermarks \cite{bui2025trustmark, sanroman2024audioseal, fernandez2024videoseal} carry a key that points to the manifest directly, while a fingerprint needs a registry lookup but survives when the watermark is removed or was never embedded. Watermarks that localize edits \cite{zhang2024editguard, zhang2025omniguard} and blind forensic detectors \cite{guillaro2023trufor} address tampering without a reference. Prior benchmarks evaluate each family on its own; we measure their interaction on the same media and attacks.

\subsection{Perceptual Quality Metrics}
Most watermarking papers report PSNR and SSIM \cite{wang2004ssim}, and some add MS-SSIM \cite{wang2003msssim} or the learned LPIPS distance \cite{zhang2018lpips}. These metrics summarize an average difference. Watermark residuals, however, are often structured patterns that a viewer notices in flat regions or when switching between the original and the marked image. FLIP models this second case: it estimates the differences an observer perceives when flipping between two images \cite{andersson2020flip}. ColorVideoVDP models spatial and temporal vision for luminance and colour and accounts for viewing conditions and display characteristics \cite{mantiuk2024cvvdp}. VMAF is a perceptual video quality metric in common use for encoder tuning \cite{rassool2017vmaf}. The Pixel Seal authors argue that proxy losses such as MSE and LPIPS fail to predict visible artefacts \cite{soucek2025pixelseal}, which is one reason to report display-aware metrics next to PSNR.

For audio, SNR and scale-invariant SNR \cite{leroux2019sdr} measure waveform error. PESQ \cite{rix2001pesq,itu2001p862} and its wideband extension \cite{itu2007p8622} predict speech quality. The ITU withdrew both recommendations in January 2024 in favour of P.863, but they are still used, which keeps results comparable with earlier watermarking papers. PESQ was designed for telephone speech and does not apply to music. STOI predicts intelligibility \cite{taal2011stoi}. Integrated loudness under ITU-R BS.1770 \cite{itu2023bs1770} and EBU R~128 \cite{ebu2023r128} detects level changes that waveform metrics do not isolate.
